\documentclass[11pt,a4paper]{article}

\usepackage[T1]{fontenc}
\usepackage[utf8]{inputenc}
\usepackage[brazil]{babel}
\usepackage{geometry}
\geometry{margin=2cm,headheight=14pt}
\usepackage{xcolor}
\usepackage{hyperref}
\usepackage{longtable}
\usepackage{array}
\usepackage{fancyhdr}

\hypersetup{
    colorlinks=true,
    linkcolor=blue!55!black,
    urlcolor=blue!55!black,
    pdftitle={Links do codigo por topico - ComposeReview},
    pdfauthor={Pedro Oliveira}
}

\pagestyle{fancy}
\fancyhf{}
\fancyhead[L]{Links do código por tópico --- ComposeReview}
\fancyhead[R]{\thepage}
\renewcommand{\headrulewidth}{0.4pt}

\newcolumntype{L}[1]{>{\raggedright\arraybackslash}p{#1}}

\newcommand{\repobase}{https://github.com/stpedr/curriculo/blob/claude/jetpack-compose-review-r09a59/jetpack-compose-review/}

\title{\textbf{Links do código por tópico} \\ \large Projeto \texttt{ComposeReview} --- tópicos fora de Jetpack Compose}
\author{Pedro Oliveira}
\date{\today}

\begin{document}
\shorthandoff{"}

\maketitle

\noindent Tabela com acesso direto, por tópico, ao código correspondente
no repositório \texttt{stpedr/curriculo}, branch
\texttt{claude/jetpack-compose-review-r09a59}. Os links levam à linha
exata onde a classe/função relevante começa, quando aplicável.

\vspace{1em}

\section*{Content Provider e compartilhamento de dados}
\begin{longtable}{L{5.4cm}L{8.5cm}}
\textbf{Arquivo} & \textbf{O que tem lá} \\ \hline \endhead
\href{\repobase app/src/main/java/com/pedro/composereview/data/NotesProvider.kt\#L22}{\url{data/NotesProvider.kt}} &
O \texttt{ContentProvider} em si: \texttt{query}/\texttt{insert}/\texttt{update}/\texttt{delete}/\texttt{getType} via \texttt{UriMatcher} \\
\href{\repobase app/src/main/java/com/pedro/composereview/data/NotesContract.kt}{\url{data/NotesContract.kt}} &
Contrato público: authority, \texttt{CONTENT\_URI}, nomes de coluna \\
\href{\repobase app/src/main/java/com/pedro/composereview/data/NotesDbHelper.kt}{\url{data/NotesDbHelper.kt}} &
\texttt{SQLiteOpenHelper} por trás do provider \\
\href{\repobase app/src/main/AndroidManifest.xml}{\url{AndroidManifest.xml}} &
Registro do \texttt{<provider>} \\
\end{longtable}

\section*{SharedPreferences e suas limitações}
\begin{longtable}{L{5.4cm}L{8.5cm}}
\textbf{Arquivo} & \textbf{O que tem lá} \\ \hline \endhead
\href{\repobase app/src/main/java/com/pedro/composereview/prefs/UserPreferences.kt\#L29}{\url{prefs/UserPreferences.kt}} &
Wrapper de \texttt{SharedPreferences} + comentário com as limitações \\
\end{longtable}

\section*{Activity/Fragment lifecycle}
\begin{longtable}{L{5.4cm}L{8.5cm}}
\textbf{Arquivo} & \textbf{O que tem lá} \\ \hline \endhead
\href{\repobase app/src/main/java/com/pedro/composereview/lifecycle/LifecycleDemoActivity.kt\#L20}{\url{lifecycle/LifecycleDemoActivity.kt}} &
\texttt{AppCompatActivity} clássica logando todo o ciclo de vida \\
\href{\repobase app/src/main/java/com/pedro/composereview/lifecycle/LifecycleDemoFragment.kt\#L26}{\url{lifecycle/LifecycleDemoFragment.kt}} &
\texttt{Fragment} com \texttt{onAttach} → \texttt{onDetach} completos \\
\href{\repobase app/src/main/res/layout/activity_lifecycle_demo.xml}{\url{res/layout/activity\_lifecycle\_demo.xml}} &
Layout XML clássico da Activity \\
\href{\repobase app/src/main/res/layout/fragment_lifecycle_demo.xml}{\url{res/layout/fragment\_lifecycle\_demo.xml}} &
Layout XML clássico do Fragment \\
\href{\repobase app/src/main/java/com/pedro/composereview/MainActivity.kt\#L59}{\url{MainActivity.kt}} &
Overrides de lifecycle da Activity principal \\
\end{longtable}

\section*{Coroutines}
\begin{longtable}{L{5.4cm}L{8.5cm}}
\textbf{Arquivo} & \textbf{O que tem lá} \\ \hline \endhead
\href{\repobase app/src/main/java/com/pedro/composereview/data/NotesRepository.kt\#L32}{\url{data/NotesRepository.kt}} &
\texttt{callbackFlow} (ContentObserver → Flow) e \texttt{withContext(Dispatchers.IO)} \\
\href{\repobase app/src/main/java/com/pedro/composereview/ui/notes/NotesViewModel.kt\#L60}{\url{ui/notes/NotesViewModel.kt}} &
\texttt{viewModelScope.launch \{ ... \}} \\
\end{longtable}

\section*{MVVM}
\begin{longtable}{L{5.4cm}L{8.5cm}}
\textbf{Arquivo} & \textbf{Camada} \\ \hline \endhead
\href{\repobase app/src/main/java/com/pedro/composereview/ui/notes/NotesViewModel.kt\#L33}{\url{ui/notes/NotesViewModel.kt}} &
ViewModel \\
\href{\repobase app/src/main/java/com/pedro/composereview/ui/notes/NotesScreen.kt\#L44}{\url{ui/notes/NotesScreen.kt}} &
View (Compose) \\
\href{\repobase app/src/main/java/com/pedro/composereview/data/NotesRepository.kt\#L23}{\url{data/NotesRepository.kt}} &
Model \\
\end{longtable}

\section*{StateFlow / LiveData}
\begin{longtable}{L{5.4cm}L{8.5cm}}
\textbf{Arquivo} & \textbf{O que tem lá} \\ \hline \endhead
\href{\repobase app/src/main/java/com/pedro/composereview/ui/notes/NotesViewModel.kt\#L46}{\url{ui/notes/NotesViewModel.kt}} &
\texttt{uiState: StateFlow<...>} (linha 46) e \texttt{noteCountLiveData: LiveData<Int>} (linha 55) \\
\href{\repobase app/src/main/java/com/pedro/composereview/ui/notes/NotesScreen.kt\#L49}{\url{ui/notes/NotesScreen.kt}} &
\texttt{collectAsStateWithLifecycle()} (linha 49) e \texttt{observeAsState()} (linha 53) \\
\end{longtable}

\vspace{1em}
\noindent\textit{Observação: os links apontam para o branch
\texttt{claude/jetpack-compose-review-r09a59}. Caso esse branch seja
mesclado ou removido, os links deixam de funcionar — nesse caso, aplique
o mesmo caminho de arquivo sobre a branch \texttt{main} do repositório.}

\end{document}
